\section{Metodologie utilizzate}
Una volta completata l'analisi prosodica ed emotiva si può passare alla parte di produzione mulsicale.

Per svolgere questo task si farà uso di un transformer

\begin{table}[H]
\centering
\caption{Esempio di tabella sintetica.}
\begin{tabular}{lcc}
\toprule
Metodo & Valore medio & Deviazione standard \\
\midrule
A & 2.34 & 0.41 \\
B & 2.85 & 0.27 \\
C & 3.12 & 0.33 \\
\bottomrule
\end{tabular}
\end{table}

\section{Stima del modello}
Un risultato di interesse \`e dato dalla seguente stima:
\begin{equation}
\hat{\beta} = \argmin_{\beta \in \R^p} \left\{ \|\vect{y} - \mat{X}\beta\|_2^2 + \lambda \|\beta\|_1 \right\}.
\end{equation}

Questo tipo di formulazione evidenzia come la regolarizzazione possa stabilire un compromesso tra accuratezza del fit e generalizzazione del modello.
